\section{Strong-Field Proposals: Saturation and Information}
\subsection{Local capacity law and its controlled consequences}
The proposed local index is
\begin{equation}
n(K)=\frac{\nu}{\nu-K},\qquad 0\leq K<\nu,
\end{equation}
with $\nu=16$. It is at least one, increases with occupancy and is unbounded as capacity is approached from below. These consequences are proved in Lean. No metric or gravitational field equation follows from this local formula alone.

The proposed clock-speed readout has a useful extension to saturation:
\begin{equation}
c_{\rm loc}(K)=c\frac{\nu-K}{\nu}.
\end{equation}
Below capacity this agrees with $c/n(K)$, is positive and does not exceed $c$. It tends continuously to zero at $K=\nu$. The divergent index is a limiting statement: Lean's division at a zero denominator is not an infinite value and must not be interpreted as a physical horizon.

The exact weak-load remainder is
\begin{equation}
n(K)=1+\frac K\nu+\frac{K^2}{\nu(\nu-K)}.
\end{equation}
A metric reconstruction is still needed to compare it with a gravitational weak-field expansion. Likewise, a local vanishing clock speed does not establish a global event horizon, which requires causal propagation in a specified spacetime.

\subsection{Inhomogeneity and the exterior-field question}
For two local loads $a,b<\nu$, Lean proves
\begin{equation}
\begin{split}
\frac{n(a)+n(b)}2-n\!\left(\frac{a+b}2\right)
=\frac{\nu(a-b)^2}{2(\nu-a)(\nu-b)(2\nu-a-b)}.
\end{split}
\end{equation}
The gap is positive for unequal loads. Thus equal mean occupancy need not give equal mean local index. This is a consequence of the nonlinear constitutive law, not a proof that two bodies of equal mass have different exterior fields. A claim of Birkhoff violation or a quantitative neutron-star deviation requires an exterior solution, matching conditions and an operational time-dilation observable; no such prediction is established here.

\subsection{Finite storage and the information question}
A saturated finite register cannot reversibly retain its old labels while encoding a fresh independent input into the same state space. Additional storage or distinguishable output outcomes are required. The finite-capacity proofs establish this constraint without specifying where an infalling state goes.

Buffering, reflection and redistribution are candidate mechanisms, but the counting obstruction does not choose one or produce a Hawking-radiation state. Information recovery, an evaporation law and a Page curve remain open. A complete account must specify the state space, its unitary update and the observables available outside the proposed saturation region.

\subsection{Entropy: the coefficient is a physical input}
For $n$ independent equal-weight cells with $q>1$ states each, the state count is $q^n$ and information in nats is $n\log q$. If each cell has area $a$, then
\begin{equation}
I=\frac{A\log q}{a},\qquad A=na.
\end{equation}
Lean proves these identities. Binary cells of area $4\ell_p^2$ give $A\log2/(4\ell_p^2)$, not $A/(4\ell_p^2)$. Choosing instead $a=4\ell_p^2\log q$ produces the target coefficient, but that choice still needs a substrate derivation. The value $D=4$ does not remove the dependence on cell states and area.

\subsection{Temperature: a conditional first-law calculation}
In units $G=\hbar=c=k_B=1$, supply the comparison geometry $A(M)=16\pi M^2$ and entropy law $S=A/4$. Then
\begin{equation}
S(M)=4\pi M^2,\qquad \frac{dS}{dM}=8\pi M.
\end{equation}
With energy $E=M$, the first law $T\,dS/dM=1$ gives
\begin{equation}
T(M)=\frac{1}{8\pi M},\qquad M>0.
\end{equation}
Lean proves the derivative, temperature relation and decrease with mass. This explains the algebraic connection between the comparison entropy and temperature coefficients. It does not derive the supplied geometry, entropy normalization or quantum radiation process by multiplying the substrate labels $\chi$, $D$ and $\pi$.
